\subsubsection{Ablation Study: Unified vs. Separate Carry Head}

To understand the contribution of architectural design choices in carry supervision, we conduct an ablation study comparing the unified head approach (Lil\={a}vati-2) against a separate carry head architecture with varying loss weighting factors. This analysis addresses a fundamental design question: should carry predictions share the same representational space as result digits, or should they be explicitly separated into distinct output heads?

\paragraph{Architectural Variants}

We compare three supervision configurations:
\begin{itemize}
\item \textbf{Lil\={a}vati-2 (Unified Head)}: Uses a single language modeling head for both result and carry positions, with equal weight ($\lambda = 1.0$ implicitly) in the unified loss computation.
\item \textbf{Separate Head ($\lambda = 0.5$)}: Uses distinct output heads for result and carry predictions, with carry loss weighted at half the result loss: $\mathcal{L} = \mathcal{L}_{\text{result}} + 0.5 \times \mathcal{L}_{\text{carry}}$.
\item \textbf{Separate Head ($\lambda = 2.0$)}: Uses distinct output heads with carry loss weighted at twice the result loss: $\mathcal{L} = \mathcal{L}_{\text{result}} + 2.0 \times \mathcal{L}_{\text{carry}}$.
\end{itemize}

The separate head architecture introduces an additional linear projection layer (carry head) that maps the same hidden state $z_H$ to carry logits, while result predictions continue to use the standard language modeling head. This architectural modification allows the model to learn specialized representations for carry propagation independently from result digit prediction.

\paragraph{Results and Analysis}

Table~\ref{tab:addition_ablation} reports accuracy metrics across all three configurations. The unified head approach (Lil\={a}vati-2) achieves near-perfect performance on both in-distribution and out-of-distribution examples, with sequence-level OOD accuracy of $99.96\%$ and carry accuracy of $99.96\%$. 

The separate head variants with different $\lambda$ values reveal important trade-offs in the supervision signal. When $\lambda = 0.5$, the separate head architecture achieves $99.53\%$ sequence-level accuracy and $99.69\%$ digit-level accuracy, with near-perfect carry accuracy ($100.00\%$). While this performance is strong, it falls short of the unified head approach, suggesting that the reduced weight on carry loss ($\lambda = 0.5$) may limit the model's ability to fully leverage carry supervision for robust generalization.

Conversely, when $\lambda = 2.0$, the increased emphasis on carry supervision results in lower sequence-level accuracy ($98.59\%$) and digit-level accuracy ($99.55\%$), despite maintaining high carry accuracy ($99.88\%$). This degradation suggests that over-weighting carry loss may cause the model to over-optimize carry prediction at the expense of result digit accuracy, or may introduce training instability due to the imbalanced loss components. The performance gap between $\lambda = 0.5$ and $\lambda = 2.0$ highlights the sensitivity of the separate head architecture to loss weighting hyperparameters.

\paragraph{Unified Head Advantage}

The superior performance of the unified head architecture (Lil\={a}vati-2) supports the hypothesis that carry and result digits share fundamental representational requirements in place-value arithmetic. By treating both as outputs of the same computational process and supervising them through a single head, the model learns a unified representation that naturally captures the relationship between result digits and their associated carry propagation. This design choice reflects the algorithmic structure of addition: carry values are not auxiliary signals but integral components of the place-value computation, and they should be represented in the same latent space as the result digits they support.

The separate head architecture, while providing explicit architectural separation, may introduce representational fragmentation that hinders the model's ability to learn the tight coupling between result and carry positions. The need to tune $\lambda$ to balance the two loss components further suggests that the separate head approach requires additional hyperparameter sensitivity that the unified approach avoids.

\paragraph{Implications for Intermediate Supervision Design}

These results suggest that when intermediate structure (like carry propagation) is algorithmically coupled with primary outputs (result digits), unified supervision through a single head is preferable to architectural separation. The unified approach eliminates the need for loss weighting hyperparameters and naturally captures the representational symmetry between related outputs. This principle may extend to other arithmetic tasks where intermediate computations are structurally related to final results, guiding future work on intermediate supervision strategies.

Figure~\ref{fig:addition_ablation} visualizes training dynamics and performance comparisons across the three configurations, highlighting the stability and convergence behavior of each approach.

\begin{table}[h]
\centering
\footnotesize
\caption{Ablation study comparing unified head (Lil\={a}vati-2) versus separate carry head architectures with different loss weighting factors $\lambda$. Results show sequence-level accuracy (Seq) and digit-level accuracy (Digit) for both in-distribution (ID) and out-of-distribution (OOD) test sets. Lil\={a}vati-2 results are from final evaluation; separate head results are from validation metrics during training. Carry Acc denotes overall carry digit accuracy.}
\label{tab:addition_ablation}
\resizebox{\linewidth}{!}{
\begin{tabular}{lcccccc}
\toprule
\textbf{Method} 
& \textbf{Head Type} 
& \textbf{$\lambda$} 
& \textbf{ID Seq} 
& \textbf{ID Digit} 
& \textbf{OOD Seq} 
& \textbf{OOD Digit} 
& \textbf{Carry Acc} \\
\midrule
Lil\={a}vati-2 
& Unified 
& 1.0 
& 1.0000 
& 1.0000 
& 0.9996 
& 1.0000 
& 0.9996 \\
Separate Head 
& Separate 
& 0.5 
& 0.9953 
& 0.9969 
& 0.9953 
& 0.9969 
& 1.0000 \\
Separate Head 
& Separate 
& 2.0 
& 0.9859 
& 0.9955 
& 0.9859 
& 0.9955 
& 0.9988 \\
\bottomrule
\end{tabular}
}
\end{table}

\begin{figure}[htbp]
\centering
\includegraphics[width=0.8\linewidth]{images/addition_ablation_comparison.png}
\caption{Training dynamics and performance comparison for unified head (Lil\={a}vati-2) versus separate carry head architectures with different loss weighting factors. Left: Validation accuracy curves over training epochs. Right: Length generalization behavior showing sequence-level accuracy as a function of result digit length.}
\label{fig:addition_ablation}
\end{figure}
